%Daten zur Institution

%\input{Dozentdaten}

%\renewcommand{\fachbereich}{Fachbereich}

%\renewcommand{\dozent}{Prof. Dr. . }

%Klausurdaten

\renewcommand{\klausurgebiet}{ }

\renewcommand{\klausurtyp}{ }

\renewcommand{\klausurdatum}{. 20}

\klausurvorspann {\fachbereich} {\klausurdatum} {\dozent} {\klausurgebiet} {\klausurtyp}

%Daten für folgende Punktetabelle

\renewcommand{\aeins}{  3  }

\renewcommand{\azwei}{  3   }

\renewcommand{\adrei}{  6  }

\renewcommand{\avier}{  4  }

\renewcommand{\afuenf}{  5  }

\renewcommand{\asechs}{  0  }

\renewcommand{\asieben}{  5  }

\renewcommand{\aacht}{  6  }

\renewcommand{\aneun}{  6  }

\renewcommand{\azehn}{  0  }

\renewcommand{\aelf}{  6  }

\renewcommand{\azwoelf}{  12  }

\renewcommand{\adreizehn}{  0   }

\renewcommand{\avierzehn}{  9  }

\renewcommand{\afuenfzehn}{  65  }

\renewcommand{\asechzehn}{    }

\renewcommand{\asiebzehn}{    }

\renewcommand{\aachtzehn}{    }

\renewcommand{\aneunzehn}{    }

\renewcommand{\azwanzig}{    }

\renewcommand{\aeinundzwanzig}{    }

\renewcommand{\azweiundzwanzig}{    }

\renewcommand{\adreiundzwanzig}{    }

\renewcommand{\avierundzwanzig}{    }

\renewcommand{\afuenfundzwanzig}{    }

\renewcommand{\asechsundzwanzig}{    }

\punktetabellevierzehn

\klausurnote

\newpage

\setcounter{section}{0}

__NOEDITSECTION__

\inputaufgabeklausurloesung
{3}

{

Definiere die folgenden
\zusatzklammer {kursiv gedruckten} {} {} Begriffe.
\aufzaehlungsechs{Eine
\stichwort {Hauptkrümmung} {}
zu einer differenzierbaren Fläche

\mavergleichskette
{\vergleichskette
{ Y
}
{ \subseteq }{ \R^3
}
{  }{
}
{    }{
}
{   }{
}
}
{}{}{}
in einem Punkt

\mavergleichskette
{\vergleichskette
{ P
}
{ \in }{ Y
}
{  }{
}
{    }{
}
{   }{
}
}
{}{}{.}

}{Ein
\stichwort {Diffeomorphismus} {}
zwischen differenzierbaren Mannigfaltigkeiten
\mathkor {} {L} {und} {M} {.}

}{Ein
\stichwort {Tangentialvektor} {}
in einem Punkt

\mavergleichskette
{\vergleichskette
{  P
}
{ \in }{  M
}
{  }{
}
{    }{
}
{   }{
}
}
{}{}{}
einer differenzierbaren Mannigfaltigkeit $M$.

}{Ein
\stichwort {Differentialoperator erster Ordnung} {}
auf einer differenzierbaren Mannigfaltigkeit $M$.

}{Ein
\stichwort {euklidischer Halbraum} {.}

}{Eine
\definitionsverweis {kompakte Ausschöpfung}{}{}
eines
\definitionsverweis {topologischen Raumes}{}{}
$X$.
}

}
{

\aufzaehlungsechs{Man nennt jeden
\definitionsverweis {Eigenwert}{}{}
der
\definitionsverweis {Weingartenabbildung}{}{}
\maabbeledisp {L_P} {T_P Y} { T_PY
} {v} {- \left(  D_{v} N  \right) \left(  P  \right)
} {,}
eine
\definitionswort {Hauptkrümmung}{}
von $Y$ in $P$.
}{Ein
\definitionsverweis {Homöomorphismus}{}{}
\maabbdisp {\varphi} {L} {M
} {}
heißt ein
\stichwortpraemath {C^1} {Diffeomorphismus}{,}
wenn sowohl
\mathkor {} {\varphi} {als auch} {\varphi^{-1}} {}
$C^1$-\definitionsverweis {Abbildungen}{}{}
sind.
}{Unter einem Tangentialvektor an $P$ versteht man eine
\definitionsverweis {Äquivalenzklasse}{}{}
von
\definitionsverweis {tangential äquivalenten}{}{}
\definitionsverweis {differenzierbaren Kurven}{}{}
durch $P$.
}{Eine Abbildung
\maabbdisp {D} {C^1(M,\R) } {C^0(M,\R)
} {}
heißt
Differentialoperator erster Ordnung,
wenn sie folgende Eigenschaften erfüllt.
\aufzaehlungzwei {$D$ ist
$\R$-\definitionsverweis {linear}{}{.}
} {Es ist

\mavergleichskette
{\vergleichskette
{ D(fg)
}
{ = }{ fD(g)+gD(f)
}
{  }{
}
{    }{
}
{   }{
}
}
{}{}{.}
}
}{Unter dem euklidischen Halbraum der Dimension $n$ versteht man die Menge

\mavergleichskettedisp
{\vergleichskette
{ H
}
{ =} { \left\{ x \in \R^n   \mid  x_1 \geq 0        \right\}
}
{ } {
}
{ } {
}
{ } {
}
}
{}{}{}
mit der
\definitionsverweis {induzierten Topologie}{}{.}
}{Eine kompakte Ausschöpfung

\mathbed {A_n} {}
 {n \in \N} {}
 {} {} {} {,}
von $X$ ist eine
\definitionsverweis {Folge}{}{}
von
\definitionsverweis {kompakten Teilmengen}{}{}

\mavergleichskette
{\vergleichskette
{ A_n
}
{ \subseteq }{ X
}
{  }{
}
{    }{
}
{   }{
}
}
{}{}{}
mit

\mathdisp {A_n \subseteq A_{n+1}^{o} \text{ und } \bigcup_{n=0}^\infty A_n = X} { . }

}

 }

__NOEDITSECTION__

\inputaufgabeklausurloesung
{3}

{

Formuliere die folgenden Sätze.
\aufzaehlungdrei{Der Satz über den
\definitionsverweis {Tangentialraum}{}{}
einer
\definitionsverweis {abgeschlossenen Untermannigfaltigkeit}{}{}

\mavergleichskette
{\vergleichskette
{  M
}
{ \subseteq }{  N
}
{  }{
}
{    }{
}
{   }{
}
}
{}{}{.}}{Der Satz über die Gestalt einer zurückgezogenen Differentialform $\varphi^* \omega$  unter einer differenzierbaren Abbildung
\maabbdisp {\varphi} { U } { V
} {}
\zusatzklammer {mit
\definitionsverweis {offenen Teilmengen}{}{}
\mathkor {} {U \subseteq \R^n} {und} {V \subseteq \R^m} {,}
deren Koordinaten mit
\mathl{x_1  , \ldots , x_n}{} bzw. mit
\mathl{y_1  , \ldots , y_m}{} bezeichnet seien} {} {}
von einer
$k$-\definitionsverweis {Differentialform}{}{}
auf $V$ mit der Darstellung

\mavergleichskettedisp
{\vergleichskette
{  \omega
}
{ =} {\sum_{I,\,  { \# \left(   I  \right) } = k}  f_I  dy_I
}
{ } {
}
{ } {
}
{ } {
}
}
{}{}{,}
wobei
\maabb {f_I} { V } {\R
} {}
\definitionsverweis {Funktionen}{}{}
sind.}{Der Satz über die Charakterisierung von geodätischen Kurven bezüglich des Levi-Civita-Zusammenhangs auf einer riemannschen Mannigfaltigkeit $M$.}

}
{

\aufzaehlungdrei{\faktsituation {Es sei $N$ eine
\definitionsverweis {differenzierbare Mannigfaltigkeit}{}{}
der
\definitionsverweis {Dimension}{}{}
$n$ und es sei

\mavergleichskette
{\vergleichskette
{ M
}
{ \subseteq }{ N
}
{  }{
}
{    }{
}
{   }{
}
}
{}{}{}
eine
\definitionsverweis {abgeschlossene Untermannigfaltigkeit}{}{}
der Dimension $m$.}
\faktfolgerung {Dann ist für jeden Punkt

\mavergleichskette
{\vergleichskette
{ P
}
{ \in }{ M
}
{  }{
}
{    }{
}
{   }{
}
}
{}{}{}
die
\definitionsverweis {Tangentialabbildung}{}{}
\maabbdisp {} { T_PM } { T_PN
} {}
\definitionsverweis {injektiv}{}{.}}
\faktzusatz {D.h. der Tangentialraum
\mathl{T_PM}{} ist ein
\definitionsverweis {Untervektorraum}{}{}
der Dimension $m$ von
\mathl{T_PN}{.}}
\faktzusatz {}}{\faktsituation {Es seien
\mathkor {} {U \subseteq \R^n} {und} {V \subseteq \R^m} {}
\definitionsverweis {offene Teilmengen}{}{,}
deren Koordinaten mit
\mathl{x_1 , \ldots , x_n}{} bzw. mit
\mathl{y_1 , \ldots , y_m}{} bezeichnet seien. Es sei
\maabbdisp {\varphi} { U } { V
} {}
eine
\definitionsverweis {differenzierbare Abbildung}{}{}
und es sei $\omega$ eine
$k$-\definitionsverweis {Differentialform}{}{}
auf $V$ mit der Darstellung

\mavergleichskettedisp
{\vergleichskette
{ \omega
}
{ =} { \sum_{I,\, { \# \left(   I  \right) } = k}  f_I  dy_I
}
{ } {
}
{ } {
}
{ } {
}
}
{}{}{,}
wobei
\maabb {f_I} { V } {\R
} {}
\definitionsverweis {
Funktionen}{}{}
sind.}
\faktfolgerung {Dann besitzt die
\definitionsverweis {zurückgezogene Form}{}{}
die Darstellung

\mavergleichskettealignhandlinks
{\vergleichskettealignhandlinks
{ \varphi^* \omega
}
{ =} { \sum_{I,\, { \# \left(   I  \right) } = k}  \left(  f_I \circ \varphi  \right)  \left( \sum_{J,\, { \# \left(   J  \right) } = k} \det  \left(  \left(  { \frac{   \partial  \varphi_i   }{  \partial   x_j   } }  \right)_{ i \in I, j \in J}  \right)  dx_J \right)
}
{ =} { \sum_{J,\, { \# \left(   J  \right) } = k}   \left( \sum_{I,\, { \# \left(   I  \right) } = k} \left(  f_I \circ \varphi  \right)  \det  \left(  \left(  { \frac{   \partial  \varphi_i   }{  \partial   x_j   } }  \right)_{ i \in I, j \in J}  \right)  \right)  dx_J
}
{ } {
}
{ } {
}
}
{}
{}{.}}
\faktzusatz {}
\faktzusatz {}}{\faktsituation {Eine zweifach
\definitionsverweis {differenzierbare Kurve}{}{}
\maabb {\gamma} { I } { M
} {}
auf einer
\definitionsverweis {riemannschen Mannigfaltigkeit}{}{}
$M$ ist}
\faktfolgerung {genau dann eine
\definitionsverweis {geodätische Kurve}{}{}
\zusatzklammer {bezüglich des
\definitionsverweis {Levi-Civita-Zusammenhangs}{}{}} {} {,}
wenn sie auf jeder Karte das
\definitionsverweis {gewöhnliche Differentialgleichungssystem zweiter Ordnung}{}{}

\mavergleichskettedisp
{\vergleichskette
{ \gamma_k^{\prime \prime} + \sum_{ij} \Gamma^k_{ij} \gamma_i' \gamma_j'
}
{ =} { 0
}
{ } {
}
{ } {
}
{ } {
}
}
{}{}{}
\zusatzklammer {für alle $k$} {} {}
erfüllt.}
\faktzusatz {}
\faktzusatz {}}

 }

__NOEDITSECTION__

\inputaufgabeklausurloesung
{6 (2+4)}

{

Es seien

\mavergleichskette
{\vergleichskette
{ \alpha,\beta
}
{ \in }{ \R_+
}
{  }{
}
{    }{
}
{   }{
}
}
{}{}{}
und sei

\mavergleichskettedisp
{\vergleichskette
{ Y
}
{ =} { \left\{ (x,y)\in\R^2  \mid   \alpha x^2+\beta y^2 = 1        \right\}
}
{ } {
}
{ } {
}
{ } {
}
}
{}{}{.}
\aufzaehlungzwei {Bestimme die
\definitionsverweis {Gauß-Abbildung}{}{}
zu $Y$.
} { Man gebe zur Gauß-Abbildung zu $Y$ explizit eine
\definitionsverweis {Umkehrabbildung}{}{}
an.
}

}
{

\aufzaehlungzwei {Als
\definitionsverweis {Einheitsnormalenfeld}{}{}
können wir vom Gradienten der beschreibenden Funktion, also von

\mavergleichskettedisp
{\vergleichskette
{ G(x,y)
}
{ =} { \begin{pmatrix}  2 \alpha x  \\ 2 \beta y    \end{pmatrix}
}
{ } {
}
{ } {
}
{ } {
}
}
{}{}{}
ausgehen, durch Normalisierung erhalten wir

\mavergleichskettedisp
{\vergleichskette
{ N(x,y)
}
{ =} {  { \frac{   1 }{  \sqrt{ \alpha^2 x^2 + \beta^2 y^2 }  } } \begin{pmatrix}   \alpha x  \\  \beta y    \end{pmatrix}
}
{ } {
}
{ } {
}
{ } {
}
}
{}{}{.}
Die Gauß-Abbildung ist also
\maabbeledisp {\varphi} {Y = \left\{  \begin{pmatrix}  x  \\ y   \end{pmatrix}   \in \R^2  \mid   \alpha x^2+\beta y^2 = 1        \right\}  } { S^1
} { \begin{pmatrix}  x  \\ y   \end{pmatrix} } { { \frac{   1 }{  \sqrt{ \alpha^2 x^2 + \beta^2 y^2 }  } } \begin{pmatrix}   \alpha x  \\  \beta y    \end{pmatrix}
} {.}
} {Wir behaupten, dass
\maabbeledisp {\psi} {S^1} {Y
} { \begin{pmatrix}  u  \\v   \end{pmatrix} } { { \frac{   1 }{  \sqrt{ { \frac{  u^2 }{ \alpha } }   + { \frac{  v^2 }{ \beta } } }  } }       \begin{pmatrix}   { \frac{   u  }{ \alpha } }   \\ { \frac{   v  }{ \beta } }    \end{pmatrix}
} {,}
eine Umkehrfunktion ist. Wegen

\mavergleichskettedisp
{\vergleichskette
{ { \frac{   1 }{   { \frac{  u^2 }{ \alpha } }   + { \frac{  v^2 }{ \beta } }  } } \left(  \alpha { \frac{  u^2 }{ \alpha^2 } } + \beta{ \frac{  v^2 }{ \beta^2 } }  \right)
}
{ =} { 1
}
{ } {
}
{ } {
}
{ } {
}
}
{}{}{}
liegt das Bild von $\psi$ in $Y$. Wegen

\mavergleichskettealign
{\vergleichskettealign
{  \varphi \left(  \psi \begin{pmatrix}  u  \\v   \end{pmatrix}  \right)
}
{ =} { \varphi \left(  { \frac{   1 }{  \sqrt{ { \frac{  u^2 }{ \alpha } }   + { \frac{  v^2 }{ \beta } } }  } }       \begin{pmatrix}   { \frac{   u  }{ \alpha } }   \\ { \frac{   v  }{ \beta } }    \end{pmatrix}    \right)
}
{ =} { \varphi \begin{pmatrix}   { \frac{   1 }{  \alpha    \sqrt{ { \frac{  u^2 }{ \alpha } }   + { \frac{  v^2 }{ \beta } } }  } }   u    \\ { \frac{   1 }{ \beta  \sqrt{ { \frac{  u^2 }{ \alpha } }  + { \frac{  v^2 }{ \beta } } }  } }  v       \end{pmatrix}
}
{ =} { { \frac{   1 }{  \sqrt{  \alpha^2 { \frac{  u^2 }{  \alpha^2   \left(   { \frac{  u^2 }{ \alpha } }  +  { \frac{  v^2 }{ \beta } }   \right)  } }  + \beta^2 { \frac{  v^2 }{  \beta^2 \left(   { \frac{  u^2 }{ \alpha } }  +  { \frac{  v^2 }{ \beta } }   \right)  } }       }  } }  \begin{pmatrix}  { \frac{   1 }{  \sqrt{ { \frac{  u^2 }{ \alpha } }   + { \frac{  v^2 }{ \beta } } }  } }   u    \\   { \frac{   1 }{   \sqrt{ { \frac{  u^2 }{ \alpha } } + { \frac{  v^2 }{ \beta } } }  } } v    \end{pmatrix}
}
{ =} { { \frac{   1 }{  \sqrt{ { \frac{  u^2 }{  \left(   { \frac{  u^2 }{ \alpha } }  +  { \frac{  v^2 }{ \beta } }   \right)  } }   + { \frac{  v^2 }{  \left(   { \frac{  u^2 }{ \alpha } }  +  { \frac{  v^2 }{ \beta } }   \right)  } }       }  } }  \begin{pmatrix}   { \frac{   1 }{  \sqrt{ { \frac{  u^2 }{ \alpha } }   + { \frac{  v^2 }{ \beta } } }  } }   u    \\   { \frac{   1 }{   \sqrt{ { \frac{  u^2 }{ \alpha } }   + { \frac{  v^2 }{ \beta } } }  } } v    \end{pmatrix}
}
}
{
\vergleichskettefortsetzungalign
{ =} { { \frac{    \sqrt{ { \frac{  u^2 }{ \alpha } }  +  { \frac{  v^2 }{ \beta } } } }{  \sqrt{    u^2+ v^2       }  } }  \begin{pmatrix}   { \frac{   1 }{  \sqrt{ { \frac{  u^2 }{ \alpha } }   + { \frac{  v^2 }{ \beta } } }  } }   u    \\   { \frac{   1 }{   \sqrt{ { \frac{  u^2 }{ \alpha } }   + { \frac{  v^2 }{ \beta } } }  } } v    \end{pmatrix}
}
{ =} { \begin{pmatrix}  u  \\v   \end{pmatrix}
}
{ } {
}
{ } {}
}
{}{}
und

\mavergleichskettealign
{\vergleichskettealign
{  \psi \left(  \varphi \begin{pmatrix}  x  \\ y   \end{pmatrix}  \right)
}
{ =} { \psi \begin{pmatrix}  { \frac{    \alpha x }{  \sqrt{\alpha^2 x^2+ \beta^2 y^2 }  } }   \\ { \frac{   \beta y }{  \sqrt{\alpha^2 x^2+ \beta^2 y^2  } }}    \end{pmatrix}
}
{ =} { { \frac{   1 }{  \sqrt{ { \frac{   \alpha x^2 }{  \alpha^2x^2 + \beta^2y^2  } }  + { \frac{   \beta y^2 }{  \alpha^2x^2 + \beta^2y^2  } }  }  } }  \begin{pmatrix}  { \frac{    x }{  \sqrt{\alpha^2 x^2+ \beta^2 y^2 }  } }   \\ { \frac{   y }{  \sqrt{\alpha^2 x^2+ \beta^2 y^2  } }}    \end{pmatrix}
}
{ =} { { \frac{   \sqrt{\alpha^2 x^2+ \beta^2 y^2 } }{  \sqrt{    \alpha x^2 +  \beta y^2  }  } }  \begin{pmatrix}  { \frac{    x }{  \sqrt{\alpha^2 x^2+ \beta^2 y^2 }  } }   \\ { \frac{   y }{  \sqrt{\alpha^2 x^2+ \beta^2 y^2  } }}    \end{pmatrix}
}
{ =} { \begin{pmatrix}  x  \\ y   \end{pmatrix}
}
}
{}
{}{}
sind die Abbildungen invers zueinander.
}

 }

__NOEDITSECTION__

\inputaufgabeklausurloesung
{4}

{

Man erläutere das Begriffspaar \anfuehrung{extrinsisch und intrinsisch}{} im Kontext von Mannigfaltigkeiten.

}
{  }

__NOEDITSECTION__

\inputaufgabeklausurloesung
{5}

{

Beweise die Formel für den Flächeninhalt der Oberfläche eines Rotationskörpers zu einer
\definitionsverweis {differenzierbaren Kurve}{}{}
\maabbeledisp {\gamma} {]a,b[} {\R^2
} { t } {(x(t),y(t))
} {,}
mit

\mavergleichskette
{\vergleichskette
{ y(t)
}
{ > }{ 0
}
{  }{
}
{    }{
}
{   }{
}
}
{}{}{.}

}
{

Es sei $S$ die Rotationsfläche, die eine abgeschlossene zweidimensionale Untermannigfaltigkeit in einer offenen Menge des $\R^3$ ist. Wir wenden
[[Flächenstück im Raum/Einbettung/Flächenform/Fakt|Korollar 17.2 (Differentialgeometrie (Osnabrück 2023))]]
auf die Parametrisierung
\maabbeledisp {} { ]a,b[ \times ]0, 2 \pi[} {S
} { (t, \alpha) } { (x(t),  y(t) \cos \alpha  , y(t) \sin \alpha  )
} {,}
an. Die partiellen Ableitungen sind

\mathdisp {\begin{pmatrix}  x'(t) \\ y'(t) \cos \alpha \\  y'(t) \sin \alpha   \end{pmatrix} \text{   und } \begin{pmatrix}  0  \\ - y(t) \sin \alpha \\  y(t) \cos \alpha   \end{pmatrix}} {  }

und daher ist

\mathdisp {E(t, \alpha) = (x'(t))^2 + (y'(t))^2,\, F(t, \alpha)=0,\, G(t, \alpha) = (y(t))^2} { . }

Somit ist der Flächeninhalt gleich

\mavergleichskettedisphandlinks
{\vergleichskettedisphandlinks
{ \int_{  a  }^{  b  } \int_{  0  }^{  2 \pi } \sqrt{(x'(t))^2+(y'(t))^2} \cdot   y(t) \, d \alpha \, d t
}
{ =} { 2 \pi \int_{  a  }^{  b  } \sqrt{(x'(t))^2+(y'(t))^2} \cdot   y(t) \, d t
}
{ } {
}
{ } {
}
{ } {
}
}
{}{}{.}

 }

__NOEDITSECTION__

\inputaufgabeklausurloesung
{0}

{

}
{  }

__NOEDITSECTION__

\inputaufgabeklausurloesung
{5}

{

Es sei $X$ ein
\definitionsverweis {Torus}{}{.}
Man gebe eine surjektive
\definitionsverweis {differenzierbare Abbildung}{}{}
\maabbdisp {\varphi} {\R^2} {X
} {}
derart an, dass auch die
\definitionsverweis {Tangentialabbildung}{}{}
\maabbdisp {T_P(\varphi)} { T_P\R^2 } { T_{\varphi(P)}X
} {}
in jedem Punkt

\mavergleichskette
{\vergleichskette
{ P
}
{ \in }{ \R^2
}
{  }{
}
{    }{
}
{   }{
}
}
{}{}{}
surjektiv ist.

}
{

Wir schreiben den Torus als $X=S^1 \times S^1$ mit dem Einheitskreis
\mathl{S^1= \left\{ (x,y) \in \R^2  \mid   \betrag { (x,y) } = 1        \right\}}{.}
Die Abbildung
\maabbeledisp {\varphi} {\R^2} {S^1 \times S^1
} {(u,v)} {( \cos  u, \sin  u , \cos  v , \sin  v )
} {,}
ist surjektiv, da es sich in den beiden Komponenten um die Standardparametrisierung des Einheitskreises handelt. Auch die Surjektivität der Tangentialabbildung lässt sich komponentenweise überprüfen, da der Tangentialraum einer Produktmannigfaltigkeit das Produkt der Tangentialräume ist. Die Ableitung von
\maabbeledisp {} {\R} { \R^2
} { t } { ( \cos  t, \sin  t)
} {,}
ist
\mathl{(- \sin  t, \cos  t) \neq (0,0)}{,} sodass dieser Bildvektor stets den eindimensionalen Tangentialraum des Einheitskreises im Bildpunkt aufspannt.

 }

__NOEDITSECTION__

\inputaufgabeklausurloesung
{6 (2+2+2)}

{

Wir betrachten den Graph $M$ der Abbildung
\maabbeledisp {\varphi} {\R^2} {\R
} {(u,v)} { u^2+uv-v^3
} {,} als zweidimensionale abgeschlossene Untermannigfaltigkeit des $\R^3$, also

\mavergleichskettedisp
{\vergleichskette
{ M
}
{ =} { \left\{ (u,v,u^2+uv-v^3)  \mid  (u,v) \in \R^2         \right\}
}
{ } {
}
{ } {
}
{ } {
}
}
{}{}{}
mit der vom $\R^3$ induzierten riemannschen Metrik. Es sei
\maabbeledisp {\psi} {\R^2} { M
} {(u,v)} { (u,v,u^2+uv-v^3)
} {,}
die zugehörige Diffeomorphie.

a) Bestimme das totale Differential zu $\psi$ sowie die Bildvektoren
\mathl{T_P(\psi) (e_1)}{} und
\mathl{T_P(\psi) (e_2)}{} in
\mathl{T_{\psi(P)}M}{.}

b) Bestimme für jeden Punkt der Form

\mavergleichskette
{\vergleichskette
{ P
}
{ = }{ (u,0)
}
{  }{
}
{    }{
}
{   }{
}
}
{}{}{}
den Flächeninhalt des von
\mathl{T_P(\psi) (e_1)}{} und
\mathl{T_P(\psi) (e_2)}{} in
\mathl{T_{\psi(P)}M}{} aufgespannten Parallelogramms.

c) Bestimme für jeden Punkt der Form

\mavergleichskette
{\vergleichskette
{ P
}
{ = }{ (0,v)
}
{  }{
}
{    }{
}
{   }{
}
}
{}{}{}
den Flächeninhalt des von
\mathl{T_P(\psi) (e_1)}{} und
\mathl{T_P(\psi) (e_2)}{} in
\mathl{T_{\psi(P)}M}{} aufgespannten Parallelogramms.

}
{

a) Das totale Differential zu $\psi$ im Punkt
\mathl{P=(u,v)}{} ist

\mathdisp {\left(D \psi\right)_{P} = \begin{pmatrix}  1 &  0 \\  0 &  1 \\   2u+v & u-3v^2 \end{pmatrix}} {  }

und es ist

\mathdisp {T_P(\psi) (e_1) = \left(  D \psi  \right)_{ P } \left(  e_1  \right) = \begin{pmatrix}  1 \\ 0\\  2u+v   \end{pmatrix} \text{   und } T_P(\psi) (e_2)   = \left(  D \psi  \right)_{ P } \left(  e_2  \right) = \begin{pmatrix}  0 \\ 1\\ u-3v^2   \end{pmatrix}} { . }

b) und c) Zur Bestimmung des Flächeninhalts berechnen wir zunächst die Skalarprodukte der beiden Vektoren
\mathkor {} {\begin{pmatrix}  1 \\ 0\\  2u+v   \end{pmatrix}} {und} {\begin{pmatrix}  0 \\ 1\\ u-3v^2   \end{pmatrix}} {.}
Es ist

\mathdisp {\left\langle \begin{pmatrix}  1 \\ 0\\  2u+v   \end{pmatrix}  , \begin{pmatrix}  1 \\ 0\\  2u+v   \end{pmatrix}   \right\rangle = 1+(2u+v)^2 = 1 +4u^2+4uv +v^2} { , }

\mathdisp {\left\langle \begin{pmatrix}  1 \\ 0\\  2u+v   \end{pmatrix}  , \begin{pmatrix}  0 \\ 1\\ u-3v^2   \end{pmatrix}    \right\rangle = (2u+v)(u-3v^2) = 2u^2 -6uv^2 +uv -3v^3} {  }

und

\mathdisp {\left\langle \begin{pmatrix}  0 \\ 1\\ u-3v^2   \end{pmatrix}   , \begin{pmatrix}  0 \\ 1\\ u-3v^2   \end{pmatrix}    \right\rangle = 1+(u-3v^2)^2 = 1 +u^2 -6uv^2  + 9v^4} { . }

b) Für
\mathl{P=(u,0)}{} berechnet sich der Flächeninhalt des von
\mathl{T_P(\psi) (e_1)}{} und
\mathl{T_P(\psi) (e_2)}{} in
\mathl{T_{\psi(P)}M}{} aufgespannten Parallelogramms zu

\mavergleichskettedisp
{\vergleichskette
{ \sqrt { \det  \begin{pmatrix}  1+4u^2  &   2u^2   \\  2 u^2 &  1+ u^2  \end{pmatrix} }
}
{ =} { \sqrt{  (1+4u^2)(1+u^2) - 4u^4 }
}
{ =} { \sqrt{  1+5u^2 }
}
{ } {
}
{ } {
}
}
{}{}{.}

c) Für
\mathl{P=(0,v)}{} berechnet sich der Flächeninhalt des von
\mathl{T_P(\psi) (e_1)}{} und
\mathl{T_P(\psi) (e_2)}{} in
\mathl{T_{\psi(P)}M}{} aufgespannten Parallelogramms zu

\mavergleichskettedisp
{\vergleichskette
{ \sqrt { \det  \begin{pmatrix}  1+v^2  &   -3v^3   \\  -3  v^3 &  1+ 9v^4  \end{pmatrix} }
}
{ =} { \sqrt{  (1+v^2)(1+9v^4) - 9v^6 }
}
{ =} { \sqrt{  1+v^2 +9v^4  }
}
{ } {
}
{ } {
}
}
{}{}{.}

 }

__NOEDITSECTION__

\inputaufgabeklausurloesung
{6}

{

Beweise die Produktregel für differenzierbare Differentialformen auf einer offenen Menge des $\R^n$.

}
{

Es seien
\mathl{x_1 , \ldots , x_n}{} die Koordinaten auf $\R^n$. Wegen der Linearität von $d$ und der
[[Dachprodukt/Algebrastruktur/Eigenschaften/Fakt|Multilinearität des Dachprodukts]]
können wir die beiden Differentialformen als
\mathkor {} {\omega=f dx_I} {und} {\tau=g dx_J} {}
mit Indexmengen
\mathkor {} {I=\{i_1 , \ldots , i_k\}} {und} {J=\{ j_1 , \ldots , j_\ell \}} {}
schreiben. Es gilt dann

\mavergleichskettealignhandlinks
{\vergleichskettealignhandlinks
{ d( \omega \wedge \tau)
}
{ =} { d \left(  fdx_I \wedge gdx_J  \right)
}
{ =} { d \left(  (f g) dx_I \wedge dx_J  \right)
}
{ =} { \sum_{s = 1}^n { \frac{   \partial  (fg)  }{  \partial   x_s  } }  dx_s  \wedge  dx_I  \wedge dx_J
}
{ =} { \sum_{s = 1}^n \left(  g { \frac{   \partial   f   }{  \partial   x_s  } } + f { \frac{   \partial   g   }{  \partial   x_s  } }  \right)  dx_s \wedge  dx_I  \wedge dx_J
}
}
{
\vergleichskettefortsetzungalign
{ =} { \sum_{s = 1}^n  g { \frac{   \partial   f   }{  \partial   x_s  } } dx_s \wedge  dx_I \wedge dx_J + \sum_{s = 1}^n  f { \frac{   \partial   g   }{  \partial   x_s  } }  dx_s \wedge  dx_I  \wedge dx_J
}
{ =} { \sum_{s = 1}^n  { \frac{   \partial   f   }{  \partial   x_s  } } dx_s \wedge  dx_I \wedge gdx_J + \sum_{s = 1}^n  { \frac{   \partial   g   }{  \partial   x_s  } }  dx_s \wedge f dx_I \wedge dx_J
}
{ =} { d(fdx_I ) \wedge g dx_J  + \sum_{s = 1}^n (-1)^k  f dx_I \wedge { \frac{   \partial   g   }{  \partial   x_s  } }  dx_s \wedge dx_J
}
{ =} { d(fdx_I ) \wedge g dx_J + (-1)^k  f dx_I \wedge  d( g dx_J)
}
}
{}{.}

 }

__NOEDITSECTION__

\inputaufgabeklausurloesung
{0}

{

}
{  }

__NOEDITSECTION__

\inputaufgabeklausurloesung
{6}

{

Es sei $M$ eine
\definitionsverweis {Mannigfaltigkeit mit Rand}{}{}
und sei

\mavergleichskette
{\vergleichskette
{ P
}
{ \in }{ \partial M
}
{  }{
}
{    }{
}
{   }{
}
}
{}{}{}
ein Randpunkt. Zeige, dass für einen Tangentialvektor

\mavergleichskette
{\vergleichskette
{ v
}
{ \in }{ T_PM
}
{  }{
}
{    }{
}
{   }{
}
}
{}{}{}
folgende Eigenschaften äquivalent sind.
\aufzaehlungzwei {Es gibt einen stetig differenzierbaren Weg
\maabb {\gamma} { [- \epsilon, 0]} { M
} {}
mit

\mavergleichskette
{\vergleichskette
{ \gamma(0)
}
{ = }{ P
}
{  }{
}
{    }{
}
{   }{
}
}
{}{}{}
und

\mavergleichskette
{\vergleichskette
{ \gamma'(0)
}
{ = }{ v
}
{  }{
}
{    }{
}
{   }{
}
}
{}{}{.}
} { $v$ wird bei jeder Karte mit einem negativen Halbraum unter der Tangentialabbildung auf den rechten Halbraum abgebildet.
}

}
{

Es sei $\gamma$ ein Halbweg in $M$ wie angegeben und sei

\mavergleichskette
{\vergleichskette
{ P
}
{ \in }{  U
}
{ \subseteq }{ M
}
{    }{
}
{   }{
}
}
{}{}{}
ein Kartengebiet mit einer Karte
\maabbdisp {\alpha} { U } {V \subseteq  H_{\leq 0}
} {}
mit

\mavergleichskettedisp
{\vergleichskette
{  H_{\leq 0}
}
{ =} { \left\{ (x_1,x_2 , \ldots, x_n)  \mid   x_1 \leq 0        \right\}
}
{ } {
}
{ } {
}
{ } {
}
}
{}{}{}
und mit

\mavergleichskettedisp
{\vergleichskette
{ Q
}
{ =} { \alpha(P)
}
{ =} { \left(  0   , \,  a_2  , \,   \ldots , \,  a_n                \right)
}
{ } {
}
{ } {
}
}
{}{}{.}
Unter der Tangentialabbildung wird

\mavergleichskette
{\vergleichskette
{ v
}
{ = }{ \gamma'(0)
}
{  }{
}
{    }{
}
{   }{
}
}
{}{}{}
auf den Ableitungsvektor von

\mavergleichskette
{\vergleichskette
{ \delta
}
{ = }{ \alpha \circ \gamma
}
{  }{
}
{    }{
}
{   }{
}
}
{}{}{}
im Nullpunkt abgebildet. Da $\gamma$ in $M$ verläuft, verläuft $\delta$ ganz in $H_{\leq 0}$. Daher ist im Differenzenquotient

\mathdisp {{ \frac{   \left(  \delta_1(t)   , \,   \delta_2(t)  , \,   \ldots , \,   \delta_n(t)                \right)  }{ t } }} {  }

sowohl $t$ negativ als auch $\delta_1(t)$ negativ und daher ist

\mavergleichskettedisp
{\vergleichskette
{  \delta_1'(0)
}
{ \geq} { 0
}
{ } {
}
{ } {
}
{ } {
}
}
{}{}{,}
gehört also zur positiven Hälfte.

Es gehöre umgekehrt
\mathl{T_P(\alpha)(v)}{} zur positiven Hälfte. Dann verläuft die Gerade
\maabbeledisp {\delta} {\R} {\R^n
} { t } { Q+tv
} {}
für

\mavergleichskette
{\vergleichskette
{ t
}
{ \leq }{  0
}
{  }{
}
{    }{
}
{   }{
}
}
{}{}{}
ganz in der negativen Hälfte. Der Weg
\mathl{\alpha^{-1} \circ \delta}{} ist auf einem negativen Intervall definiert und landet in $M$ und ist eine Halbwegrealisierung von $v$.

 }

__NOEDITSECTION__

\inputaufgabeklausurloesung
{12}

{

Beweise den Satz von Stokes für Mannigfaltigkeiten mit Rand.

}
{

\teilbeweis {}{}{}
{Es sei

\mathbed {U_i} {}
 {i \in I} {}
 {} {} {} {,}
eine
\definitionsverweis {offene Überdeckung}{}{}
von $M$ mit
\definitionsverweis {orientierten Karten}{}{}
und es sei

\mathbed {h_j} {}
 {j \in J} {}
 {} {} {} {,}
eine dieser
\definitionsverweis {Überdeckung untergeordnete stetig differenzierbare Partition der Eins}{}{,}
die nach
[[Mannigfaltigkeit/Abzählbare Basis/Überdeckung/Partition/Fakt|Satz 22.10 (Differentialgeometrie (Osnabrück 2023))]]
existiert. Zu jedem

\mavergleichskette
{\vergleichskette
{ P
}
{ \in }{ M
}
{  }{
}
{    }{
}
{   }{
}
}
{}{}{}
gibt es eine offene Umgebung

\mavergleichskette
{\vergleichskette
{ P
}
{ \in }{ V_P
}
{ \subseteq }{ M
}
{    }{
}
{   }{
}
}
{}{}{}
derart, dass

\mavergleichskette
{\vergleichskette
{ h_j \vert_{V_P }
}
{ = }{ 0
}
{  }{
}
{    }{
}
{   }{
}
}
{}{}{}
bis auf endlich viele Ausnahmen ist. Es sei $Y$ der Träger von $\omega$. Die Überdeckung

\mavergleichskette
{\vergleichskette
{ Y
}
{ \subseteq }{ \bigcup_{P \in M} V_P
}
{  }{
}
{    }{
}
{   }{
}
}
{}{}{}
besitzt wegen der vorausgesetzten
\definitionsverweis {Kompaktheit}{}{}
eine endliche Teilüberdeckung, sagen wir

\mavergleichskettedisp
{\vergleichskette
{ Y
}
{ \subseteq} { V_{1} \cup \ldots \cup V_{r}
}
{ =} { V
}
{ } {
}
{ } {
}
}
{}{}{.}
Daher sind überhaupt nur endlich viele der $h_j$ auf $V$ von $0$ verschieden. Wir setzen

\mavergleichskette
{\vergleichskette
{ \omega_j
}
{ = }{ h_j \omega
}
{  }{
}
{    }{
}
{   }{
}
}
{}{}{;}
diese Differentialformen sind ebenfalls stetig differenzierbar. Der Träger von $h_j$ ist eine in $M$
\definitionsverweis {abgeschlossene Teilmenge}{}{,}
die in
\mathl{U_{i(j)}}{} liegt, daher liegt der Träger von $\omega_j$ in
\mathl{Y \cap U_{i(j)}}{} und ist selbst kompakt nach
[[Kompakter Raum/Abgeschlossene Teilmenge/Kompakt/Fakt/Beweis/Aufgabe|Aufgabe 13.10 (Differentialgeometrie (Osnabrück 2023))]].
Es gilt

\mavergleichskettedisp
{\vergleichskette
{ \omega
}
{ =} { \sum_{j \in J} \omega_j
}
{ } {
}
{ } {
}
{ } {
}
}
{}{}{,}
wobei nur endlich viele dieser Differentialformen
\mathl{\omega_j}{} von $0$ verschieden sind, da

\mavergleichskette
{\vergleichskette
{ \omega_j \vert_{M \setminus Y}
}
{ = }{ 0
}
{  }{
}
{    }{
}
{   }{
}
}
{}{}{}
für alle $j$ ist und

\mavergleichskettedisp
{\vergleichskette
{ \omega_j \vert_V
}
{ =} { 0
}
{ } {
}
{ } {
}
{ } {
}
}
{}{}{}
für alle $j$ bis auf endlich viele Ausnahmen. Wegen der
[[Volumenform/Integration/Eigenschaften/Fakt|Additivität des Integrals von Differentialformen]]
und der
[[Differentialform/Mannigfaltigkeit/Äußere Ableitung/Eigenschaften/Fakt|Additivität der äußeren Ableitung]]
kann man die Aussage für die einzelnen $\omega_j$ getrennt beweisen. Wir können also annehmen, dass eine stetig differenzierbare $(n-1)$-Differentialform gegeben ist, die kompakten Träger besitzt, der ganz in einer Kartenumgebung $U$ liegt.}
{}
\teilbeweis {}{}{}
{Es liegt ein
\definitionsverweis {Diffeomorphismus}{}{}
\maabb {\alpha} { U } { V
} {}
mit

\mavergleichskette
{\vergleichskette
{ V
}
{ \subseteq }{ H
}
{  }{
}
{    }{
}
{   }{
}
}
{}{}{}
offen vor, der zugleich einen Diffeomorphismus zwischen den Rändern

\mavergleichskette
{\vergleichskette
{ \partial U
}
{ = }{ \partial M \cap U
}
{  }{
}
{    }{
}
{   }{
}
}
{}{}{}
und

\mavergleichskette
{\vergleichskette
{ \partial V
}
{ = }{ \partial H \cap V
}
{  }{
}
{    }{
}
{   }{
}
}
{}{}{}
induziert. Dabei gilt

\mavergleichskettedisp
{\vergleichskette
{
\int_{  \partial U  }   \omega
}
{ =} {
\int_{  \partial V  }  \alpha^{-1 *}\omega
}
{ } {
}
{ } {
}
{ } {
}
}
{}{}{}
und

\mavergleichskettedisp
{\vergleichskette
{
\int_{ U  }  d \omega
}
{ =} {
\int_{ V  }  \alpha^{-1 *} (  d \omega)
}
{ =} {
\int_{  V  }  d \left(  \alpha^{-1 *}\omega  \right)
}
{ } {
}
{ } {
}
}
{}{}{}
nach
[[Differentialform/Offene Menge/Äußere Ableitung/Eigenschaften/Fakt|Lemma 20.2 (Differentialgeometrie (Osnabrück 2023))  (5)]].}
{\leerzeichen{}Wir können also von einer auf

\mavergleichskette
{\vergleichskette
{ V
}
{ \subseteq }{ H
}
{  }{
}
{    }{
}
{   }{
}
}
{}{}{}
definierten stetig differenzierbaren Differentialform mit kompaktem Träger ausgehen, die wir auf ganz $H$ außerhalb des Trägers als Nullform fortsetzen können.}
\teilbeweis {}{}{}
{Wegen der Kompaktheit des Trägers gibt es einen hinreichend großen Quader

\mavergleichskette
{\vergleichskette
{ Q
}
{ \subseteq }{ H
}
{  }{
}
{    }{
}
{   }{
}
}
{}{}{,}
dessen eine Seite $S$ auf $\partial H$ liegt und der den Träger von $\omega$ nur in $S$ trifft. Auf allen anderen Seiten von $Q$ ist $\omega$
\zusatzklammer {und damit auch $d\omega$} {} {}
die Nullform. Daher gilt einerseits

\mavergleichskettedisp
{\vergleichskette
{
\int_{ H  }  d\omega
}
{ =} {
\int_{ Q  }  d\omega
}
{ } {
}
{ } {
}
{ } {
}
}
{}{}{}
und andererseits

\mavergleichskettedisp
{\vergleichskette
{
\int_{  \partial H  }   \omega
}
{ =} {
\int_{  S  }   \omega
}
{ =} {
\int_{  \partial Q  }   \omega
}
{ } {
}
{ } {
}
}
{}{}{.}
Somit folgt die Aussage
aus der [[Satz von Stokes/Quaderversion/Fakt|Quaderversion des Satzes von Stokes]].}
{}

 }

__NOEDITSECTION__

\inputaufgabeklausurloesung
{0}

{

}
{  }

__NOEDITSECTION__

\inputaufgabeklausurloesung
{9 (1+1+3+2+2)}

{

Wir betrachten die trigonometrische Parametrisierung
\maabbeledisp {\psi} {\R } { S^1
} { t } { \left(  \cos  t   , \,   \sin  t                   \right)
} {}
des Einheitskreises. Es sei

\mavergleichskette
{\vergleichskette
{ c
}
{ \in }{ [0, 2 \pi[
}
{  }{
}
{    }{
}
{   }{
}
}
{}{}{}
fixiert. Auf dem trivialen Vektorbündel
\maabbdisp {} {S^1 \times \R^2} { S^1
} {}
sei ein
\definitionsverweis {linearer Zusammenhang}{}{}
$\nabla$ gegeben derart, dass der längs $\psi$
\definitionsverweis {zurückgezogene Zusammenhang}{}{}
$\psi^* \nabla$ durch die
\definitionsverweis {Christoffelsymbole}{}{}
\zusatzklammer {der Index für die einzige Ableitungsrichtung wird weggelassen} {} {}

\mavergleichskettedisp
{\vergleichskette
{ \begin{pmatrix} \tilde{\Gamma}_1^1 (t)   &  \tilde{\Gamma}_1^2 (t)   \\ \tilde{\Gamma}_2^1(t)  &  \tilde{\Gamma}_2^2(t)  \end{pmatrix}
}
{ =} { \begin{pmatrix}  0  &   - { \frac{  c  }{  2 \pi   } }   \\  { \frac{  c  }{  2 \pi  } }  &  0  \end{pmatrix}
}
{ } {
}
{ } {
}
{ } {
}
}
{}{}{}
gegeben ist. Wir betrachten zu einem Punkt

\mavergleichskette
{\vergleichskette
{ Q
}
{ = }{ \begin{pmatrix}  a  \\b   \end{pmatrix}
}
{ \in }{ \R^2
}
{    }{
}
{   }{
}
}
{}{}{}
die Abbildung
\maabbeledisp {\Psi_{Q}} {\R} { S^1 \times \R^2
} { t } { \left(  \cos  t   , \,   \sin  t  , \,   M(t)  \begin{pmatrix}  a  \\b   \end{pmatrix}                 \right)
} {}
mit

\mavergleichskettedisp
{\vergleichskette
{ M(t)
}
{ =} { \begin{pmatrix}  \cos  { \frac{  ct  }{  2 \pi   } }   &   -  \sin  { \frac{  ct  }{  2 \pi   } }    \\  \sin  { \frac{  ct  }{  2 \pi  } }   &  \cos  { \frac{  ct  }{  2 \pi   } }  \end{pmatrix}
}
{ } {
}
{ } {
}
{ } {
}
}
{}{}{.}
\aufzaehlungfuenf{Bestimme $\Psi_Q(0)$.
}{Bestimme $\Psi_Q(2 \pi)$.
}{Zeige, dass  $\Psi_Q$ eine
\definitionsverweis {horizontale Liftung}{}{}
längs $\psi$  ist.
}{Zeige, dass
\maabbeledisp {} {\R^2} { \R^2
} { Q } { \Psi_Q(2 \pi)
} {,}
eine
\definitionsverweis {lineare Isometrie}{}{}
ist
\zusatzklammer {der Basispunkt $(1,0)$ wird hier nicht aufgeführt} {} {.}
}{Zeige, dass
\maabbeledisp {} {\Z} { \operatorname{SL}_{  2  } \! \left(  \R  \right)
} { k } { \left(  Q \mapsto \Psi_Q(2k \pi )  \right)
} {,}
ein
\definitionsverweis {Gruppenhomomorphismus}{}{}
ist.
}

}
{

\aufzaehlungfuenf{Es ist

\mavergleichskettealign
{\vergleichskettealign
{ \Psi_Q(0)
}
{ =} { \left(  \cos  0   , \,   \sin  0  , \,   M(0) \begin{pmatrix}  a  \\b   \end{pmatrix}                 \right)
}
{ =} { \left(  1   , \,   0  , \,   \begin{pmatrix}   1   &   0   \\  0   &  1  \end{pmatrix}    \begin{pmatrix}  a  \\b   \end{pmatrix}                 \right)
}
{ =} { \left(  1   , \,   0  , \,     \begin{pmatrix}  a  \\b   \end{pmatrix}                 \right)
}
{ } {
}
}
{}
{}{.}
}{Es ist

\mavergleichskettealign
{\vergleichskettealign
{ \Psi_Q(2 \pi )
}
{ =} { \left(  \cos  2 \pi   , \,   \sin  2 \pi  , \,   M(2 \pi) \begin{pmatrix}  a  \\b   \end{pmatrix}                 \right)
}
{ =} { \left(  1   , \,   0  , \,   \begin{pmatrix}  \cos  c   &   - \sin c    \\  \sin c   &   \cos  c     \end{pmatrix} \begin{pmatrix}  a  \\b   \end{pmatrix}                 \right)
}
{ =} { \left(  1   , \,   0  , \,     \begin{pmatrix} a \cos  c   - b \sin c   \\ a \sin  c   - b \cos c    \end{pmatrix}                 \right)
}
{ } {
}
}
{}
{}{.}
}{Es ist

\mavergleichskettealign
{\vergleichskettealign
{  \Psi_Q(t)
}
{ =} { \left(  \cos  t    , \,   \sin  t  , \,   \begin{pmatrix}   \cos  { \frac{  ct  }{  2 \pi   } }   &   - \sin  { \frac{  ct  }{  2 \pi   } }    \\  \sin  { \frac{  ct  }{  2 \pi   } }   &   \cos  { \frac{  ct  }{  2 \pi   } }  \end{pmatrix}  \begin{pmatrix}  a  \\b   \end{pmatrix}                 \right)
}
{ =} { \left(  \cos  t   , \,   \sin  t  , \,   \begin{pmatrix}  a \cos  { \frac{  ct  }{  2 \pi  } }  - b \sin  { \frac{  ct  }{  2 \pi  } }   \\ a \sin  { \frac{  ct  }{  2 \pi  } }  +   b  \cos  { \frac{  ct  }{  2 \pi } }    \end{pmatrix}                 \right)
}
{ =} { \left(  \cos  t   , \,   \sin  t  , \,     \left(  a \cos  { \frac{  ct  }{  2 \pi  } }  - b \sin  { \frac{  ct  }{  2 \pi  } }  \right)   e_1 +  \left(  a \sin  { \frac{  ct  }{  2 \pi  } }  +   b  \cos  { \frac{  ct  }{  2 \pi } }  \right)   e_2                 \right)
}
{ } {
}
}
{}
{}{,}
insbesondere liegt eine Liftung vor. Die vertikale Ableitung dieses Schnittes in $\R^2$ über $\R$ kann man mit den angegebenen Christoffelsymbolen bestimmen, es ist

\mavergleichskettealigndrucklinks
{\vergleichskettealigndrucklinks
{ \left(  \psi^* \nabla  \right)_\partial  \left(    \left(  a \cos  { \frac{  ct  }{  2 \pi  } }  - b \sin  { \frac{  ct  }{  2 \pi  } }  \right)   e_1 +  \left(  a \sin  { \frac{  ct  }{  2 \pi  } }  +   b  \cos  { \frac{  ct  }{  2 \pi } }  \right)   e_2  \right)
}
{ =} {   \left(  a \cos  { \frac{  ct  }{  2 \pi  } }  - b \sin  { \frac{  ct  }{  2 \pi  } }  \right)   \left(  \psi^* \nabla  \right)_\partial \left(  e_1  \right)  + \partial    \left(  a \cos  { \frac{  ct  }{  2 \pi  } }  - b \sin  { \frac{  ct  }{  2 \pi  } }  \right)   e_1    +   \left(  a \sin  { \frac{  ct  }{  2 \pi  } }  +   b  \cos  { \frac{  ct  }{  2 \pi } }  \right)      \left(  \psi^* \nabla  \right)_\partial \left(  e_2  \right) + \partial  \left(  a \sin  { \frac{  ct  }{  2 \pi  } }  +   b  \cos  { \frac{  ct  }{  2 \pi } }  \right)   e_2
}
{ =} { -  \left(  a \cos  { \frac{  ct  }{  2 \pi  } }  - b \sin  { \frac{  ct  }{  2 \pi  } }  \right)  { \frac{   c  }{  2 \pi } }  e_2  +   { \frac{   c  }{  2 \pi } } \left(  -a \sin  { \frac{  ct  }{  2 \pi  } }  - b \cos  { \frac{  ct  }{  2 \pi  } }  \right)   e_1    +   \left(  a \sin  { \frac{  ct  }{  2 \pi  } }  +   b  \cos  { \frac{  ct  }{  2 \pi } }  \right) { \frac{   c  }{  2 \pi } }   e_1    +  { \frac{   c  }{  2 \pi } }   \left(  a \cos  { \frac{  ct  }{  2 \pi  } } -  b  \sin  { \frac{  ct  }{  2 \pi } }  \right)   e_2
}
{ =} { 0
}
{ } {
}
}
{}{}{,}
also liegt eine
\definitionsverweis {horizontale Liftung}{}{}
vor.
}{Es geht um die Abbildung
\maabbeledisp {} {\R^2} { \R^2
} { \begin{pmatrix}  a  \\b   \end{pmatrix} } {    \begin{pmatrix}  \cos  c   &   - \sin  c   \\  \sin  c  &  \cos  c  \end{pmatrix} \begin{pmatrix}  a  \\b   \end{pmatrix}
} {,}
die linear ist. Die beschreibende Matrix ist orthogonal
\zusatzklammer {eine Drehmatrix} {} {,}
deshalb liegt eine Isometrie vor.
}{Es ist
\mathl{\Psi (2k \pi)}{} die lineare Abbildung
\maabbeledisp {} {\R^2} { \R^2
} {\begin{pmatrix}  a  \\b   \end{pmatrix} } { \begin{pmatrix}  \cos kc   &   - \sin kc   \\  \sin kc  &  \cos kc  \end{pmatrix} \begin{pmatrix}  a  \\b   \end{pmatrix}
} {,}
also die Drehung um den Winkel $kc$. Die Drehmatrix zur Summe von zwei Winkeln ist das Matrixprodukt der beiden Drehmatrizen
\zusatzklammer {siehe
[[Sinus und Kosinus/Reell/Additionstheoreme/Drehung/Fakt|Fakt]] [[Kurs:Differentialgeometrie/Sinus und Kosinus/Reell/Additionstheoreme/Drehung/Fakt/Faktreferenznummer|*****]]} {} {,}
daher liegt ein
\definitionsverweis {Gruppenhomomorphismus}{}{}
vor.
}

 }

 [[Kategorie:Latexseite]]
